% 孔径问题教学构造：边缘方向(1,1)，真实运动(2,0)，约束线u-v=2；图像坐标纵轴向下。宽约100mm。
\begin{tikzpicture}[x=1mm,y=1mm,font=\figurefont\small,
  vec/.style={-{Stealth[length=2mm]},line width=1.1pt},
  ax/.style={-{Stealth[length=1.8mm]},draw=FigureStroke,line width=.7pt},
  lab/.style={font=\figurefont\footnotesize,text=FigureMuted}]
  % 左：小孔中的边缘
  \begin{scope}
    \clip (0,0) circle (16);
    \fill[FigureBlueFill] (-20,20) -- (20,-20) -- (20,20) -- cycle;
    \draw[FigureStroke,line width=.8pt] (-20,20) -- (20,-20);
    \draw[FigureMuted,dashed,line width=.6pt] (-8,20) -- (32,-20);
  \end{scope}
  \draw[FigureStroke,line width=.8pt] (0,0) circle (16);
  \draw[vec,FigureRed] (0,0) -- (12,0);
  \draw[vec,FigureRed,dashed] (0,0) -- (0,12);
  \draw[vec,FigureBlue] (0,0) -- (6,6);
  \node[lab] at (0,-19.5) {小孔内只能看到一段直边缘};
  % 右：速度平面，v轴向下
  \begin{scope}[shift={(58,0)}]
    \draw[ax] (-6,0) -- (18,0) node[right] {$u$};
    \draw[ax] (0,14) -- (0,-14) node[below] {$v$};
    \draw[FigureStroke,line width=1pt] (-4,16) -- (16,-4);
    \fill[FigureRed] (12,0) circle (1.1);
    \fill[FigureRed] (0,12) circle (1.1);
    \fill[FigureBlue] (6,6) circle (1.1);
    \draw[FigureMuted,dashed,line width=.6pt] (0,0) -- (6,6);
    \node[anchor=south west,lab] at (12.6,0.6) {$(2,0)$};
    \node[anchor=east,lab] at (-1,12) {$(0,-2)$};
    \node[anchor=west,lab] at (7,7.5) {法向光流$(1,-1)$};
    \node[anchor=west,lab] at (3,-8) {约束线$u-v=2$};
  \end{scope}
\end{tikzpicture}
